
\begin{proof}[引理~\ref{lem:7-4-64} 的证明]
除符号外，结论显然。
符号可按顶点数归纳核对；
也可依据图~\ref{fig:7-4-5}
附近的讨论核对。
\end{proof}

现在利用引理~\ref{lem:7-4-64}
完成命题~\ref{prop:7-4-56}
的证明。

\par\noindent\textbf{命题~\ref{prop:7-4-56} 的证明。}
先注意到
$\mathfrak g(\mathbf Y)$
的像包含于 $\overline H[1]$，
因此
\begin{equation}
0=
\overline{\mathfrak m}_1
\bigl(\mathfrak g(\mathbf Y)(\mathbf x)\bigr)
-(-1)^{\deg'\mathbf Y+1}
\mathfrak g(\mathbf Y)
\bigl(
\widehat{\overline{\mathfrak m}}_1(\mathbf x)
\bigr)
\tag{7.4.69}\label{eq:7-4-69}
\end{equation}
对
$\mathbf x\in B(\overline H[1])
\subset B(\overline C[1])$
成立。
注意
$\deg\mathfrak g(\mathbf Y)
=\deg'\mathbf Y+1$。
我们对
$(D,\delta)
=(\overline C[1],\overline{\mathfrak m}_1)$
应用引理~\ref{lem:7-4-64}，
计算式~\eqref{eq:7-4-69}
右端。
先暂不计符号，
稍后再加以讨论。

由于
$\mathfrak h(|\Gamma|,
Z_\Gamma^{\mathfrak g})
=\pm\mathfrak g(\Gamma)$，
引理~\ref{lem:7-4-64}
表明式~\eqref{eq:7-4-69}
右端等于
\begin{equation}
\begin{aligned}
\sum_{\Gamma\in G(k,\mathbf Y)}
T^{E(\Gamma)}e^{\mu(\Gamma)/2}
\Bigl(
&\sum_e\pm
\mathfrak h\bigl(
|\Gamma|,\partial_e
Z_\Gamma^{\mathfrak g}
\bigr)(\mathbf x)\\
&+
\sum_v\pm
\mathfrak h\bigl(
|\Gamma|,\partial_v
Z_\Gamma^{\mathfrak g}
\bigr)(\mathbf x)
\Bigr).
\end{aligned}
\tag{7.4.70}\label{eq:7-4-70}
\end{equation}

先看式~\eqref{eq:7-4-70}
的第一项。
若 $e$ 为外边，则
$Z_\Gamma^{\mathfrak g}(e)
=\operatorname{id}$ 或 $\Pi_H$，
故与这些外边对应的求和项为零。
若 $e$ 为内边，则
$Z_\Gamma^{\mathfrak g}(e)
=\mathfrak G$，
于是
\[
\partial_e
Z_\Gamma^{\mathfrak g}(e)
=\overline{\mathfrak m}_1\circ\mathfrak G
+\mathfrak G\circ\overline{\mathfrak m}_1
=-(\operatorname{id}-\Pi_H).
\]
据此引入两套装饰：
\[
\begin{aligned}
\partial_e Z_\Gamma^{\mathfrak g,I}(f)
&=
\begin{cases}
Z_\Gamma^{\mathfrak g}(f),
&f\ne e,\\
\operatorname{id},
&f=e,
\end{cases}
\\[0.6ex]
\partial_e Z_\Gamma^{\mathfrak g,II}(f)
&=
\begin{cases}
Z_\Gamma^{\mathfrak g}(f),
&f\ne e,\\
\Pi_H,
&f=e.
\end{cases}
\end{aligned}
\]
由定义，
\[
\begin{aligned}
\partial_e Z_\Gamma^{\mathfrak g}(e)
&=-\bigl(
\partial_e Z_\Gamma^{\mathfrak g,I}(e)
-\partial_e Z_\Gamma^{\mathfrak g,II}(e)
\bigr)\\
&=
\partial_e Z_\Gamma^{\mathfrak g,II}(e)
-\partial_e Z_\Gamma^{\mathfrak g,I}(e).
\end{aligned}
\]

如果边 $e$ 的两个端点都是
$M$ 型或都是 $Y$ 型，
则分别称 $e$ 为 $M$ 型或 $Y$ 型；
若二者均非，则称为 $MY$ 型。
分别记所有这三类内边的集合为
$C_M^1(\Gamma)$、
$C_Y^1(\Gamma)$
和 $C_{MY}^1(\Gamma)$。
于是三类内边恰构成内边集\translatornote{源书写 $C^1(\Gamma)=C_M^1(\Gamma)\cup C_Y^1(\Gamma)$；右端既漏去刚定义的 $MY$ 型内边，也不含外边。译文将等式限定于内边集并补齐第三类。见疑误 A-0575。}：
\[
C_{\mathrm{int}}^1(\Gamma)
=C_M^1(\Gamma)
\cup C_Y^1(\Gamma)
\cup C_{MY}^1(\Gamma).
\]

现在有如下引理。

\begin{lemma}{7.4.71}\label{lem:7-4-71}
\[
\begin{aligned}
&\sum_{\Gamma\in G(k,\mathbf Y)}
T^{E(\Gamma)}e^{\mu(\Gamma)/2}
\Bigl(
\sum_{e\in C_M^1(\Gamma)}
\pm\mathfrak h\bigl(
|\Gamma|,
\partial_e Z_\Gamma^{\mathfrak g,I}
\bigr)\\
&\hspace{36mm}
+
\sum_{v\in C_M^0(\Gamma)}
\pm\mathfrak h\bigl(
|\Gamma|,
\partial_v Z_\Gamma^{\mathfrak g}
\bigr)
\Bigr)=0.
\end{aligned}
\]
\end{lemma}

\begin{proof}
第一项由如下复合构成：
\[
\sum\pm
\mathfrak m_{c_1,\beta_1}
\circ
\widehat{\mathfrak m}_{c_2,\beta_2},
\]
并与其他算子复合。
这里对
$c_1,c_2,\beta_1,\beta_2$
求和，在给定 $a,\beta$ 时满足
$c_1+c_2=a+1$、
$\beta_1+\beta_2=\beta$，
且
$(c_i,\beta_i)\ne(1,0)$。
见图~\ref{fig:7-4-7}。
\par\medskip\noindent
\begin{minipage}{0.49\textwidth}
\centering
\setlength{\unitlength}{0.65mm}
\begin{picture}(87,69)
\linethickness{0.55pt}
\put(44,56){\line(0,-1){11}}
\multiput(44,58)(0,2){5}{\line(0,1){0.8}}
\put(44,45){\line(0,-1){22}}
\qbezier(44,45)(28,37)(12,29)
\qbezier(44,45)(38,37)(32,29)
\qbezier(44,45)(48,37)(52,29)
\qbezier(44,45)(56,37)(68,29)
\qbezier(44,23)(30,15)(16,7)
\qbezier(44,23)(37,15)(31,7)
\qbezier(44,23)(48,15)(52,7)
\qbezier(44,23)(58,15)(72,7)
\multiput(29,5)(-1,-2){3}{\line(0,-1){0.8}}
\multiput(54,5)(1,-2){3}{\line(0,-1){0.8}}
\put(47,48){$\scriptstyle\mathfrak m_{c_1,\beta_1}$}
\put(45,34){$\scriptstyle\operatorname{id}$}
\put(46,22){$\scriptstyle\mathfrak m_{c_2,\beta_2}$}
\end{picture}
\par\smallskip
{\renewcommand{\thefigure}{7.4.7}\refstepcounter{figure}
\normalfont\bfseries 图~\thefigure\label{fig:7-4-7}}
\end{minipage}\hfill
\begin{minipage}{0.49\textwidth}
\centering
\setlength{\unitlength}{0.65mm}
\begin{picture}(87,69)
\linethickness{0.55pt}
\put(44,56){\line(0,-1){11}}
\multiput(44,58)(0,2){5}{\line(0,1){0.8}}
\put(44,45){\line(0,-1){22}}
\qbezier(44,45)(28,37)(12,29)
\qbezier(44,45)(38,37)(32,29)
\qbezier(44,45)(48,37)(52,29)
\qbezier(44,45)(56,37)(68,29)
\qbezier(44,23)(30,15)(16,7)
\qbezier(44,23)(37,15)(31,7)
\qbezier(44,23)(48,15)(52,7)
\qbezier(44,23)(58,15)(72,7)
\multiput(29,5)(-1,-2){3}{\line(0,-1){0.8}}
\multiput(54,5)(1,-2){3}{\line(0,-1){0.8}}
\put(47,48){$\scriptstyle Y_i$}
\put(45,34){$\scriptstyle\operatorname{id}$}
\put(46,22){$\scriptstyle Y_j$}
\end{picture}
\par\smallskip
{\renewcommand{\thefigure}{7.4.8}\refstepcounter{figure}
\normalfont\bfseries 图~\thefigure\label{fig:7-4-8}}
\end{minipage}\par\medskip


另一方面，
由式~\eqref{eq:7-4-63}，
第二项由如下复合构成：
\[
\pm\mathfrak m_{a,\beta}
\circ\widehat{\overline{\mathfrak m}}_1
\ \pm\
\overline{\mathfrak m}_1
\circ\widehat{\mathfrak m}_{a,\beta},
\]
并与其他算子复合。
由 $A_\infty$ 公式，
这两部分互相抵消。
\end{proof}

置
\[
\mathbf Y_{[i,j]}
=\bigl(
[Y_i,Y_j],
Y_1,\ldots,\widehat Y_i,\ldots,
\widehat Y_j,\ldots,Y_\ell
\bigr).
\]
则不难看出
\begin{equation}
\begin{aligned}
&\sum_{\Gamma\in G(k,\mathbf Y)}
\sum_{e\in C_Y^1(\Gamma)}
\pm T^{E(\Gamma)}e^{\mu(\Gamma)/2}
\mathfrak h\bigl(
|\Gamma|,
\partial_e Z_\Gamma^{\mathfrak g,I}
\bigr)\\
&\qquad=
\sum_{i<j}
\sum_{\Gamma\in G(k,\mathbf Y_{[i,j]})}
\pm T^{E(\Gamma)}e^{\mu(\Gamma)/2}
\mathfrak h\bigl(
|\Gamma|,Z_\Gamma^{\mathfrak g}
\bigr).
\end{aligned}
\tag{7.4.72}\label{eq:7-4-72}
\end{equation}
见图~\ref{fig:7-4-8}。

再置
\[
\mathbf Y_i
=\bigl(
Y_1,\ldots,\delta(Y_i),\ldots,Y_\ell
\bigr),
\]
其中 $\delta$ 为 Hochschild 微分。

\begin{lemma}{7.4.73}\label{lem:7-4-73}
\[
\begin{aligned}
&\sum_i
\sum_{\Gamma\in G(k,\mathbf Y_i)}
\pm T^{E(\Gamma)}e^{\mu(\Gamma)/2}
\mathfrak h\bigl(
|\Gamma|,Z_\Gamma^{\mathfrak g}
\bigr)\\
&\qquad=
\sum_{\Gamma\in G(k,\mathbf Y)}
T^{E(\Gamma)}e^{\mu(\Gamma)/2}
\Bigl(
\sum_{e\in C_{MY}^1(\Gamma)}
\pm\mathfrak h\bigl(
|\Gamma|,
\partial_e Z_\Gamma^{\mathfrak g,I}
\bigr)\\
&\hspace{51mm}
+
\sum_{v\in C_Y^0(\Gamma)}
\pm\mathfrak h\bigl(
|\Gamma|,
\partial_v Z_\Gamma^{\mathfrak g}
\bigr)
\Bigr).
\end{aligned}
\]
\end{lemma}
